\section{总结与延伸}

\subsection{本讲知识体系}

本讲把”训练一个模型”拆成了可操作的最小单元：
\begin{itemize}
    \item 张量如何占内存——从 dtype、numel 到 GPU 内存层次；
    \item 矩阵乘法如何计算 FLOPs——从 $2BDK$ 到 $6ND$ 规则；
    \item 参数如何封装进 \texttt{nn.Module}——从 \texttt{nn.Parameter} 到 \texttt{state\_dict}；
    \item 数据如何切 batch——memmap、pinned memory、异步拷贝；
    \item 优化器、训练循环、checkpoint 和混合精度如何协同工作。
\end{itemize}

\subsection{核心公式速查}

\begin{table}[H]
\centering
\begin{tabular}{p{0.28\textwidth}p{0.32\textwidth}p{0.28\textwidth}}
\toprule
\textbf{场景} & \textbf{公式} & \textbf{用途} \\
\midrule
张量内存 & $\text{numel} \times \text{element\_size}$ & 单个 tensor 内存估算 \\
矩阵乘法 FLOPs & $2BDK$ & 单次 matmul 计算量 \\
训练总 FLOPs & $6ND$ & 完整训练计算预算 \\
训练时间 & $C / (\text{GPUs} \times \eta \times \text{Peak})$ & Wall-clock 估算 \\
AdamW 静态内存 & $16N$ bytes & 混合精度下每参数 \\
Transformer 每层参数 & $\approx 12d^2$ & 模型规模速算 \\
Activation checkpointing & $\sqrt{L} \cdot a$ 内存，+33\% 计算 & 内存优化评估 \\
\bottomrule
\end{tabular}
\caption{本讲核心公式速查表}
\end{table}

\subsection{从本讲到后续课程的衔接}

\begin{table}[H]
\centering
\begin{tabular}{p{0.22\textwidth}p{0.32\textwidth}p{0.30\textwidth}}
\toprule
\textbf{本讲内容} & \textbf{后续延伸} & \textbf{对应课程} \\
\midrule
内存核算 & 模型并行、张量并行、流水线并行 & Lecture 5 (GPUs) \\
FLOPs 估算 & Scaling laws、compute-optimal 训练 & Lecture 4 (Scaling) \\
\texttt{nn.Module} & 完整 Transformer 实现 & Assignment 1 \\
混合精度 & 分布式混合精度、FSDP & Lecture 5 (GPUs) \\
Checkpoint & 分布式 checkpoint、容错训练 & 高级话题 \\
\bottomrule
\end{tabular}
\caption{本讲内容在课程体系中的位置}
\end{table}

\begin{importantbox}{最后要记住的一句话}
深度学习系统不是”把模型写出来”就结束了，而是要持续回答三个问题：
\begin{enumerate}
    \item 这个东西占多少内存？
    \item 这个东西消耗多少 FLOPs？
    \item 这个东西能不能稳定地、可复现地训练起来？
\end{enumerate}
\end{importantbox}

做完 Assignment 1 之后，这些概念会变得更具体，也会更牢固，因为你会亲手把它们用在真正的 Transformer 上。

\subsection{本章小结}
这一讲把 PyTorch 中真正决定训练是否可行的三件事放到了一起：参数与状态的内存核算、前向反向的 FLOPs 规模，以及 mixed precision / checkpoint 等工程优化。真正有价值的 takeaway 不是某个公式，而是把“模型实现”提升成“系统预算”的习惯。

\subsection{拓展阅读}

\begin{itemize}
    \item \textbf{混合精度训练原始论文}：Micikevicius et al., ``Mixed Precision Training'' (2018). 提出了 loss scaling 和 master weights 的混合精度框架。
    \item \textbf{AdamW 论文}：Loshchilov \& Hutter, ``Decoupled Weight Decay Regularization'' (2019). 阐明了解耦权重衰减与 L2 正则化的区别。
    \item \textbf{Activation Checkpointing 原始论文}：Chen et al., ``Training Deep Nets with Sublinear Memory Cost'' (2016). 提出了 $\sqrt{L}$ 策略的理论分析。
    \item \textbf{FP8 训练}：Micikevicius et al., ``FP8 Formats for Deep Learning'' (2022). NVIDIA 提出的 E4M3/E5M2 FP8 标准。
    \item \textbf{PyTorch 内部机制}：Edward Z. Yang 的博客 ``PyTorch internals''. 详细讲解了 tensor storage、autograd 和 dispatch 机制。
    \item \textbf{Scaling Laws}：Kaplan et al., ``Scaling Laws for Neural Language Models'' (2020). 建立了参数量、数据量、计算量与 loss 之间的幂律关系。
    \item \textbf{Chinchilla}：Hoffmann et al., ``Training Compute-Optimal Large Language Models'' (2022). 给出了更精确的 compute-optimal 训练配方。
\end{itemize}

\end{document}
